\chapter{Catatonia}

\section*{Definition}
Catatonia is a psychomotor syndrome involving characteristic changes in movement, behavior, and responsiveness, ranging from stupor and mutism to agitation and echophenomena. Diagnosis is clinical and requires a characteristic cluster of signs; a rating scale supports assessment but does not by itself establish the diagnosis.

\section*{When to See a Doctor}

\begin{itemize}
 \item Immobility, mutism, stupor, or bizarre posturing develops acutely
 \item There is extreme agitation with purposeless excitement
 \item Autonomic instability, fever, or rigidity suggests malignant catatonia or neuroleptic malignant syndrome
\end{itemize}

\section*{How It Develops}
Catatonia results from dysregulation of GABAergic and glutamatergic signaling in cortical--basal ganglia--thalamocortical motor circuits. Reduced GABA-A activity in the orbitofrontal and parietal cortex disrupts motor inhibition, while NMDA receptor hypofunction may contribute to the motor and behavioral abnormalities.

\section*{Causes}
\begin{itemize}
 \item \textbf{Psychiatric:} Bipolar disorder (manic or depressive episodes), schizophrenia, schizoaffective disorder, major depressive disorder
 \item \textbf{Neurological:} Anti-NMDA receptor encephalitis, stroke, Parkinson disease, epilepsy (non-convulsive status), brain tumors, traumatic brain injury
 \item \textbf{Medical:} Metabolic disturbances (hyponatremia, hepatic encephalopathy), infections (encephalitis, meningitis), autoimmune disorders, drug withdrawal
 \item \textbf{Toxic:} Neuroleptic malignant syndrome, serotonin syndrome, cocaine or phencyclidine intoxication
\end{itemize}

\section*{Related Symptoms}
\begin{itemize}
 \item Mutism, stupor, waxy flexibility, posturing, stereotypy, echolalia, echopraxia, grimacing, negativism, excitement
\end{itemize}
\textbf{Important:} Catatonia is underdiagnosed in psychiatric settings and can be life-threatening, especially in the malignant form.

\section*{Medical Evaluation}
\begin{itemize}
 \item Clinical observation and the Bush-Francis Catatonia Rating Scale (BFCRS)
\\item Rule out medical and neurological causes with targeted tests such as electrolytes, glucose, calcium, liver and renal function, thyroid testing, and toxicology when indicated
 \item Neuroimaging (MRI brain) and EEG if seizure or encephalitis suspected
 \item Lumbar puncture for suspected autoimmune encephalitis (anti-NMDA receptor antibodies)
\end{itemize}

\section*{Treatment Options}
\begin{itemize}
 \item First-line treatment is usually a clinician-supervised benzodiazepine trial, often lorazepam; dosing, route, monitoring, and response vary by patient
 \item Second-line: Electroconvulsive therapy (ECT), especially for malignant catatonia
 \item If neuroleptic malignant syndrome or serotonin syndrome is suspected, stop the causative medicine under urgent medical supervision and provide syndrome-specific supportive treatment
 \item NMDA receptor encephalitis: immunotherapy (steroids, IVIG, plasmapheresis)
\end{itemize}

\section*{Self-Care and Home Management}
\begin{itemize}
 \item Catatonia requires urgent medical evaluation; home management is not appropriate for acute episodes
 \item Provide a quiet, low-stimulation environment if symptoms are mild
 \item Ensure safety and prevent falls or injury
\end{itemize}

\section*{References}
\begin{thebibliography}{99}
\bibitem{catatonia:ref1} Fink M, Taylor MA. ``Catatonia: A Clinician's Guide to Diagnosis and Treatment.'' Cambridge University Press, 2003.
\bibitem{catatonia:ref2} Bush G, Fink M, Petrides G, et al. ``Catatonia: I. Rating Scale and Standardized Examination.'' Acta Psychiatr Scand. 1996;93(2):129--136.
\bibitem{catatonia:ref3} Rosebush PI, Mazurek MF. ``Catatonia and Its Treatment.'' Schizophr Bull. 2010;36(2):239--242.
\bibitem{catatonia:ref4} Luchini F, Medda P, Rizzato S, et al. ``Electroconvulsive Therapy in Catatonia: A Systematic Review.'' J ECT. 2015;31(4):225--232.
\bibitem{catatonia:ref5} Sienaert P, Dhossche DM, Vancampfort D, et al. ``A Clinical Review of the Treatment of Catatonia.'' Front Psychiatry. 2014;5:181.
\bibitem{catatonia:ref6} American Psychiatric Association. ``Diagnostic and Statistical Manual of Mental Disorders.'' 5th ed, Text Revision. APA, 2022.
\end{thebibliography}
